\input pdfToolbox
\input pdfmsym
\pdfmsymsetscalefactor{10}
\input ../preamble
\def\enumindent{0cm}

\def\printmcount{\the\counter{section}.\the\counter{math counter}}
\setcounter{section}{10}

\centerline{\setfontandscale{bf}{20pt}Commutative Algebra \& Algebraic Geometry}
\medskip
\centerline{\bf Exercise \the\counter{section}}
\smallskip
\centerline{\it Ari Feiglin}

\def\pgl{{\rm PGL}}
\def\aut{{\rm Aut}}

\bprob

    Let $n\in\bN$, and denote by $\aut(\bP^n)$ the group of automorphisms of $\bP^n$.
    \benum
        \item Construct a homomorphism $\eta\colon\gl_{n+1}(K)\to\aut(\bP^n)$.
        \item Prove that the center of $\gl_{n+1}(K)$ consists of scalar matrices $K^\times I$.
        \item Deduce that $\eta$ factors as $\gl_{n+1}(K)\to\pgl_{n+1}(K)\xto{\bar\eta}\aut(\bP^n)$ where $\bar\eta$ is injective.
        \item Let $\c H$ be the set of hyperplanes of $\bP^n$.
        Show that $\aut(\bP^n)$ acts transitively on $\c H$.
        \item Let $P_0,\dots,P_k\in\bP^n$ be distinct points.
        Show that there is a hyperplane $H\in\c H$ such that $P_0\in H$ but $P_i\notin H$ for all other $i$.
    \eenum

\eprob

\benum
    \item Let $g\in\gl_{n+1}(K)$, then $g$ defines a map $\bA^{n+1}-0\to\bA^{n+1}-0$ as it is injective.
    Since it is linear it acts on its coordinates linearly, thus defining a morphism.
    By problem 3 in exercise 9, this means it defines a unique morphism $\eta(g)\colon\bP^n\to\bP^n$ such that $\pi\circ g=\eta(g)\circ\pi$.
    Then $\eta$ is a homomorphism: for $g,h\in\gl_{n+1}(K)$,
    $$ \eta(g\circ h)\circ\pi = \pi\circ(g\circ h) = (\pi\circ g)\circ h = \eta(g)\circ \pi\circ h = \eta(g)\circ\eta(h)\circ\pi . $$
    By uniqueness of $\eta(g)$, this means that $\eta(g\circ h)=\eta(g)\circ\eta(h)$.
    Since $\eta(\id)=\id$, this means that $\eta$ maps $\gl_{n+1}(K)$ to automorphisms of $\bP^n$ homomorphically.

    \item If $g\in Z(\gl_{n+1}(K))$, then $g$ commutes with elementary matrices.
    In particular, for $c\in K^\times$ let $c\rho_i$ denote the diagonal matrix with ones on the diagonal except for $c$ at the $i$th index.
    Then $(c\rho_i)g$ corresponds to the elementary row operation of multiplying $g$'s $i$th row by $c$, and $g(c\rho_i)$ is multiplying $g$'s $i$th column by $c$.
    Since $(c\rho_i)g=g(c\rho_i)$ this clearly means that $g$ must be a diagonal matrix.

    We claim that $g$ is in fact a scalar matrix: so $g_{ii}=g_{jj}$ for all $i,j$.
    Let $\rho_{ij}$ be the elementary matrix for swapping the $i$th and $j$th rows/columns.
    So $\rho_{ij}g$ is the matrix obtained by swapping $g$'s $i$th and $j$th rows, and $g\rho_{ij}$ is swapping $g$'s $i$th and $j$th columns.
    Then $\rho_{ij}g$'s element at the $(i,j)$th index is $g_{jj}$, while $g\rho_{ij}$'s is $g_{ii}$.
    Since $\rho_{ij}g=g\rho_{ij}$, this means $g_{jj}=g_{ii}$ as desired.

    So all matrices in the center are scalar.
    The converse is clear, since multiplying by a scalar matrix (on either side) is the same as scaling the matrix.

    \item We have $\eta(g)=\id$ iff $\pi\circ g=\pi$, by uniqueness.
    This means that for every $P\in\bA^{n+1}-0$, $g(P)$ is a scalar multiple of $P$ (since $\pi\circ g(P)=\pi(P)$).
    So all vectors are eigenvectors of $g$, which means that $g$ must be scalar.
    Conversely a scalar matrix $g$ clearly satisfies $\pi\circ g=\pi$.

    Thus $\ker\eta=Z(\gl_{n+1}(K))$, and so $\eta$ factors as desired.

    \item For a hyperplane $H=Z\parens{\sum_ia_ix_i}$, $H$ is the set of all $\pi(b_0,\dots,b_n)$ such that $\sum_ia_ib_i=0$.
    For $g\in\gl_{n+1}(K)$, let us conflate $g$ and $\eta(g)$ so we can write $g\circ\pi=\pi\circ g$.
    Then
    $$ \pi(\bar b) \in gH \iff g^{-1}\pi(\bar b) \in H \iff \pi(g^{-1}\bar b) \in H \iff \bar a\cdot(g^{-1}\bar b) = 0 , $$
    and this in turn is iff $((g^\top)^{-1}\bar a)\cdot\bar b=0$, which is iff $\bar b$ is in the hyperplane given by coefficients $(g^\top)^{-1}\bar a$.
    So $\gl_{n+1}(K)$ action restricts to one on $\c H$.
    Furthermore since $\gl_{n+1}(K)$ acts transitively on $\bA^{n+1}-0$, it acts transitively on $\c H$: for any $\bar c$ we can find $g\in\gl_{n+1}(K)$ such that $(g^\top)^{-1}\bar a=\bar c$.

    \item First let us assume $P_0=[1:0:\cdots:0]$.
    Then a hyperplane $H$ containing $P_0$ is any one of the form $H=Z\parens{\sum_{i>0}a_ix_i}$ (i.e.\ $a_0=0$).
    Let $P_i=[b_{i,0}:\cdots:b_{i,n}]$; since $P_i\neq P_0$ we have that $b_{i,j}\neq0$ for some $j>0$.
    A hyperplane containing both $P_0$ and $P_i$ is of the form
    $$ H = Z\parens{\sum_{j>0}a_jx_j},\qquad \sum_{i>0}a_jb_{i,j} = 0 . $$

    So there is a one-to-one correspondence between hyperplanes containing $P_0$ and $P_i$ and vectors in the kernel of $\phi_i\colon(a_1,\dots,a_n)\mapsto\sum_{i=1}^na_jb_{i,j}$.
    Since the $b_{i,j}$ are not all zero, this kernel is not all of $\bA^n$.
    Thus $\bigcup_i\ker\phi_i$ is not all of $\bA^n$ (since then it would be a vector space, and the finite union of vector spaces is a vector space iff one contains all the others, which would mean
    $\ker\phi_i=\bA^n$ for some $i$).

    So let $(a_1,\dots,a_n)$ not in $\bigcup_i\ker\phi_i$, then the hyperplane it defines $H=Z\parens{\sum_{i=1}^na_ix_i}$ contains $P_0$ but not contain any other $P_i$.

    In general, since $\gl_{n+1}(K)$ acts transitively on $\bA^{n+1}-0$ and thus $\bP^n$, there is a $g$ such that $gP_0=[1:0:\cdots:0]$.
    Then let $H$ be such that $gP_0\in H$ but $gP_i\notin H$ for all other $i$.
    Then $P_0\in g^{-1}H$ but $P_i\notin g^{-1}H$ for all other $i$, so $g^{-1}H$ gives the desired hyperplane.
\eenum

\bprob

    Let $n\geq1$.
    \benum
        \item Prove that restriction forms an isomorphism $\O(\bA^{n+1})\to\O(\bA^{n+1}-0)$.
        \item Deduce that $\O(\bP^n)\cong k$ and thus that $\bP^n$ is not an affine variety.
    \eenum

\eprob

\benum
    \item First we show that restriction is injective.
    Suppose that $f|_{\bA^{n+1}-0}=0$ for $f\in k[\bA^{n+1}]$.
    Thus $Z(f)\supseteq\bA^{n+1}-0$ which is open, and thus dense since $\bA^{n+1}$ is irreducible, meaning $Z(f)=\bA^{n+1}$.
    Thus $f=0$ as desired.

    Since $\bA^{n+1}-0$ is open, its coordinate ring consists of rational functions $f/g$ where $g$ is nonzero on $\bA^{n+1}-0$.
    If $g$ is nonzero on $\bA^{n+1}-0$, this means that $Z(g)\subseteq\set0$.
    But if $g$ is not a scalar, then $Z(g)$ has dimension $n-1$ and thus cannot be empty or a point; a contradiction.
    Thus $g$ is a scalar, and $f/g$ is a polynomial.

    So sections of $\bA^{n+1}-0$ are polynomials, showing the surjectivity of restriction.

    \item By definition
    $$ \O(\bP^n) = \set{f\colon\bP^n\to k}[{f\circ\pi\in\O(\bA^{n+1}-0)=k[\bA^{n+1}]}] . $$
    Suppose $f\circ\pi=p\in k[\bA^{n+1}]$, then this means for all $\lambda\in k^\times$
    $$ p(\lambda\bar x) = f\circ\pi(\lambda\bar x) = f\circ\pi(\bar x) = p(\bar x) . $$
    Writing $p=\sum_\alpha a_\alpha\bar x^\alpha$ this means that
    $$ \sum_\alpha \lambda^{\abs\alpha}a_\alpha x^\alpha = \sum_\alpha a_\alpha x^\alpha . $$
    Viewing this as a polynomial in $\lambda$, this means all of its non-constant coefficients are zero:
    $$ \sum_{\abs\alpha=n}a_\alpha x^\alpha = 0 $$
    for all $n>0$.
    So for $\abs\alpha>0$ we have $a_\alpha=0$.
    And thus $p$ is a scalar.

    Thus $f\circ\pi$ is a constant scalar, and so $f$ is as well (by surjectivity of $\pi$).
    Thus $\O(\bP^n)=k$ as desired.

    If $\bP^n$ were affine, its coordinate ring would be a quotient of $k[\bA^m]$ for some $m$, by an ideal $I$ and $\bP^n\cong Z(I)$.
    But since its coordinate ring is $k$, $I$ would have to be maximal, and then $Z(I)$ is a point.
    $\bP^n$ is not a point, so it cannot be affine.
\eenum

\bprob

    Prove that
    \benum
        \item being an isomorphism is local on the target,
        \item being an isomorphism is not local on the source.
    \eenum

\eprob

\benum
    \item First we claim that if $\phi\colon X\to Y$ is a morphism of SWFs and $U\subseteq Y$ is open, then the restriction
    $\phi\colon\phi^{-1}U\to U$ is a morphism of SWFs as well.
    This is clear since we are just restricting our perspective to open subsets of $U$ and $\phi^{-1}U$, which are open subsets of $Y$
    and $X$ respectively.

    Now if $\phi\colon X\to Y$ is an isomorphism, let $U\subseteq Y$ be open.
    Then the restriction $\phi\colon\phi^{-1}U\to U$ remains a morphism with inverse $\phi^{-1}\colon U\to\phi^{-1}U$ (since this is just
    the restriction of the morphism $\phi^{-1}$).
    Thus $\phi|_{\phi^{-1}U}$ is an isomorphism; in particular if $\phi$ is an isomorphism then it is locally an isomorphism on its target.

    Conversely suppose $\set{U_i}$ is an open cover of $Y$ such that $\phi\colon\phi^{-1}U_i\to U_i$ is an isomorphism for each $i$.
    Then $\phi$ is injective since if $\phi x=\phi y$, then these images lie in some $U_i$, so $x,y\in\phi^{-1}U_i$ over which $\phi$
    is bijective and thus $x=y$.
    And $\phi$ is surjective, as every element of $U_i$ has a preimage in $\phi^{-1}U_i$ since $\phi$ is an isomorphism there, and the
    $U_i$ cover $Y$.

    Thus $\phi^{-1}$ exists globally, and it is continuous since restricting it to the open cover $\phi^{-1}U_i$ is continuous.
    It is also a morphism, since taking a basis of $X$ by open subsets of $\phi^{-1}U_i$ as $i$ ranges, we see that $\phi^{-1}$ defines
    a morphism of coordinate rings from these basic open subsets, since $\phi^{-1}$ is a morphism $U_i\to\phi^{-1}U_i$.
    By the equivalence of SWFs defined on basic open subsets and SWFs defined on all open subsets, we have the desired result.

    \item Consider the morphism $\phi=(\id,\id)\colon\bA^n\amalg\bA^n\to\bA^n$.
    Restricting $\phi$ to each copy of $\bA^n$ in the source yields $\id$, which is an isomorphism, but clearly $\phi$ itself is not
    as it is not even injective.
\eenum

\bprob

    Prove that:
    \benum
        \item a composition of affine morphisms is affine,
        \item an isomorphism of varieties is affine,
        \item being an affine morphism is local on the target,
        \item being an affine morphism is not local on the source.
    \eenum

\eprob

\benum
    \item Suppose $X\xtos\phi Y\xtos\psi Z$ are affine.
    Let $U\subseteq Z$ be open affine, then $\psi^{-1}U\subseteq Y$ is open affine as $\psi$ is an affine morphism, and thus
    $(\psi\circ\phi)^{-1}U=\phi^{-1}(\psi^{-1}U)$ is affine as $\phi$ is affine.
    Thus $\psi\circ\phi$ is affine.

    \item Suppose $X\xtos\phi Y$ is an isomorphism, and let $U\subseteq Y$ be open affine.
    Then $\phi\colon\phi^{-1}U\to U$ is an isomorphism (being an isomorphism is local on the source), so $\phi^{-1}U$ is isomorphic
    to an affine variety (as $U$ is), and thus affine.
    Thus $\phi$ is affine.

    \item Clearly if $\phi$ is affine, it is locally affine on the target; we just take a covering of the target by open affine subsets $U_i$,
    then $\phi^{-1}U_i$ is affine by assumption.
    In recitation we showed that $\phi$ is affine iff there is an open cover of affine subsets $\set{U_i}$, $\phi^{-1}U_i$ is affine.
    So conversely suppose there is an open cover of the target $\set{W_i}$ such that $\phi$ is affine $\phi^{-1}W_i\to W_i$.

    Since open subsets of algebraic varieties are algebraic varieties, we can cover each $W_i$ by an open covering of affine subsets
    $\set{U_i^j}$, and then by assumption $\phi|_{\phi^{-1}W_i}^{-1}(U_i^j)=\phi^{-1}(U_i^j)$ is affine for each $i,j$.
    Thus there is an open affine covering $\set{U_i^j}$ for which $\phi^{-1}U_i^j$ is affine for all $i,j$.
    By the result from recitation, this means that $\phi$ is affine.

    \item Consider $\phi=(\id,\id)\colon\bA^1\amalg\bA^1\to\bA^1$.
    This is locally affine on its source, as it is locally an isomorphism on its source and isomorphisms are affine.
    But if $\phi$ were affine, this would imply that $\phi^{-1}\bA^1=\bA^1\amalg\bA^1$ is affine.
    The disjoint union has coordinate ring
    $$ \O(\bA^1\amalg\bA^1) \cong \O(\bA^1)\times\O(\bA^1) \cong k[x]\times k[y] . $$
    Define $f\colon k[x,y]\to k[x]\times k[y]$ in the obvious manner, mapping $x$ to $(x,0)$ and $y$ to $(0,y)$.
    This then maps
    $$ \sum_{n,k}a_{n,k}x^ny^k \mapsto \sum_{n,k}a_{n,k}(x^n,0)(0,y^k) = \sum_na_{n,0}(x^n,0) + \sum_ka_{0,k}(0,y^k) . $$
    So its kernel is $(xy)$.

    Since these are affine varieties, this means $\bA^1\amalg\bA^1\cong Z(xy)\subseteq\bA^2$.
    The irreducible components of $Z(xy)$ are $Z(x),Z(y)$ (they are irreducible since their ideals are prime, and their union is $Z(xy)$).
    Since both copies of $\bA^1$ are the irreducible components of $\bA^1\amalg\bA^1$, this means that the isomorphism must map
    one copy to $Z(x)$ and the other to $Z(y)$.
    But $Z(x)\cap Z(y)=Z(x,y)=\set0$ is nonempty, so the isomorphism is not injective, a contradiction.

    So $\phi$ is not affine, despite being locally affine on its source.
\eenum

\advancecounter{math counter}{1}

\bprob

    Let $X$ be a topological space, show the following.
    \benum
        \item If $Y\subseteq X$ is a subspace, then $\dim Y\leq\dim X$.
        \item If $X=\bigcup_iX_i$ is a finite union of closed subsets, then $\dim X=\max_i\dim X_i$.
        \item If $X=\bigcup_iU_i$ is an open cover, $\dim X=\sup_i\dim U_i$.
        \item Give an example of $X$ and an open dense $U\subseteq X$ such that $\dim U<\dim X$.
        \item Give an example of an infinite-dimensional Noetherian space.
    \eenum

\eprob

\benum
    \item Suppose $Y_0\subseteq Y$ is closed irreducible, then we claim $\cll_X(Y_0)\subseteq X$ is closed irreducible.
    Indeed, suppose $\cll_X(Y_0)=Z_0\cup Z_1$ are closed sets, then
    $$ Y_0 = \cll_X(Y_0)\cap Y = (Z_0\cap Y) \cup (Z_1\cap Y) , $$
    where the first equality is because $\cll_Y(Y_0)=\cll_X(Y_0)\cap Y$ and $Y_0$ is closed.
    Since $Y_0$ is irreducible and $Z_i\cap Y$ are closed in $Y$, $Y_0=Z_i\cap Y$ for some $i$.
    Thus
    $$ \cll_X(Y_0) = \cll_X(Z_i\cap Y) \subseteq \cll_X(Z_i) = Z_i , $$
    which means $\cll_X(Y_0)=Z_i$, so $\cll_X(Y_0)$ is indeed irreducible.

    Now we claim that if $Y_0\subsetneq Y_1\subseteq Y$ are closed, then $\cll_X(Y_0)\subsetneq\cll_X(Y_1)$.
    Weak inclusion is clear.
    For strict inclusion let $y\in Y_1-Y_0$, then $y\in\cll_X(Y_1)$ but it cannot be in $\cll_X(Y_0)$.
    If it were, then $y\in\cll_X(Y_0)\cap Y=\cll_Y(Y_0)=Y_0$ a contradiction.

    Thus if $\emptyset\neq Y_0\subsetneq\cdots\subsetneq Y_n$ is a chain of irreducible closed subsets of $Y$, then $\emptyset\neq\cl{Y_0}\subsetneq\cdots\subsetneq\cl{Y_n}$
    is a chain of irreducible closed subsets of $X$.
    Thus $\dim Y\leq\dim X$ as desired.

    \item By point (1), we have $\dim X_i\leq\dim X$ for all $i$ and thus $\max_i\dim X_i\leq\dim X$.
    Now let $\emptyset\neq Z_0\subsetneq\cdots\subsetneq Z_n$ be a chain of irreducible closed subsets of $X$.
    We have
    $$ Z_n = \bigcup_i(Z_n\cap X_i) $$
    and since $Z_n$ is irreducible and this is a finite union of closed sets, there is some $i$ where $Z_n=Z_n\cap X_i$, meaning $Z_n\subseteq X_i$.
    Thus this irreducible chain is contained in $X_i$; thus $\dim X\leq\max_i\dim X_i$ as desired.

    \item Again by point (1) we have $\sup_i\dim U_i\leq\dim X$.
    Now let $\emptyset\neq Z_0\subsetneq\cdots\subsetneq Z_n$ be a chain of irreducible closed subsets of $X$.
    Let $U_i$ be such that $Z_0\cap U_i$ is nonempty (since $Z_0$ is nonempty and the $U_i$ cover $X$, such an $i$ exists).
    Since $Z_j\cap U_i$ is open in $Z_j$ which is irreducible, and thus $Z_j\cap U_i$ is irreducible (open subsets of irreducible spaces are irreducible because irreducibility is equivalent
    to non-empty open subsets having non-empty intersection, if this is true for a space $X$ it is clearly true for all open subsets of $X$).

    Thus $\emptyset\neq Z_0\cap U_i\subseteq\cdots\subseteq Z_n\cap U_i$ forms a chain of irreducible subsets of $U_i$, and they are all closed since the $Z_j$ are.
    Furthermore, we claim $Z_j\cap U_i\subsetneq Z_{j+1}\cap U_i$.
    This is because $Z_{j+1}\cap U_i$ is dense in $Z_{j+1}$ (nonempty open subsets of irreducible spaces are dense), but $Z_j\cap U_i$ is not: it is dense in $Z_j$ but $Z_j$ is closed and thus
    not dense in $Z_{j+1}$.
    So this is a proper chain of closed irreducible subsets of $U_i$, meaning $\dim X\leq\sup_i\dim U_i$ as desired.

    \item Let $X=\set{0,1}$ with the topology $\set{\emptyset,X,\set0}$; the closed subsets of $X$ are then $\set{\emptyset,X,\set1}$.
    Then $\set 0\subseteq X$ is open dense as the only closed subset it is contained in is $X$.
    But as a point, $\dim\set0=0$.

    On the other hand $X$ is irreducible, since the only proper closed subset of $X$ is $\set1$ and it is not equal to $X$.
    Thus $\emptyset\neq\set1\subsetneq X$ is a chain of closed irreducible subsets of $X$, so $\dim X=1$.
    So $\dim\set0<\dim X$ despite $\set0$ being open dense in $X$.

    \item Consider $\omega$ (the ordinal) equipped with the topology whose closed sets are $[n]=\set{0,\dots,n-1}$ for $n\leq\omega$ (setting $[0]=\emptyset$ and $[\omega]=\omega$).
    This is a topology as
    $$ [n] \cup [m] = [\max\set{n,m}],\qquad \bigcap_i[n_i] = [\min_in_i] . $$
    Furthermore this shows that each $[n]$ is irreducible, since if $[n]=[n_0]\cup[n_1]$ then $n=\max\set{n_0,n_1}$, so $n=n_0$ or $n=n_1$.
    Thus we have arbitrarily long chains of irreducible closed subsets
    $$ \emptyset \neq [1] \subsetneq [2] \subsetneq \cdots \subsetneq [n] , $$
    thus $\dim\omega=\infty$.

    And $\omega$ is Noetherian: let $\Sigma$ be a collection of closed subsets of $\omega$.
    Let $\Sigma_\omega$ be the set of $n\leq\omega$ for which $[n]\in\Sigma$.
    Since $\omega$ is well-ordered, $\Sigma_\omega$ has a minimum, which means $[\min\Sigma_\omega]\in\Sigma$.
    Then
    $$ \bigcap\Sigma = \bigcap_{[n]\in\Sigma}[n] = [\min\set{n\leq\omega}[{[n]\in\Sigma}] = [\min\Sigma_\omega] \in \Sigma . $$
    Thus $\Sigma$ has a minimal element; so $\omega$'s closed subsets of the DCC and $\omega$ is Noetherian.
\eenum

\bye

